\documentclass[11pt,a4paper]{article}
\usepackage[review]{acl}
% \usepackage{acl}  % uncomment for camera-ready (remove [review])
\usepackage{times,latexsym}
\usepackage{amsmath,amssymb}
\usepackage{graphicx}
\graphicspath{{../../figures/paper/}}
\usepackage{booktabs}
\usepackage{multirow}
\usepackage{xcolor}
\usepackage{algorithm}
\usepackage{algorithmic}

\title{Less Is More: Exception Events Outperform Complete Traces for Verdict Prediction}
\author{Anonymous ACL submission}

\begin{document}
\maketitle

\begin{abstract}
Code language models increasingly incorporate execution traces, yet more trace information does not guarantee better performance.
We present the first systematic component ablation of trace representations for verdict prediction on LLM-generated Python code---classifying whether a test passes or fails---across four models spanning three architecture families and a 1.5B--13B parameter range.
Exception events alone achieve ${\geq}$99.33\% accuracy across all four models and all seeds tested under \texttt{exception\_only} under \texttt{exception\_only}, while complete traces reach only ${\sim}$80\% despite containing strictly more information.
Training dynamics reveal why: \texttt{exception\_only} converges within 50 training steps, whereas even \texttt{full\_trace} with label-only loss (gradient dilution eliminated) oscillates throughout training without reaching stable accuracy---overfitting to spurious patterns in verbose content rather than learning the verdict signal.
Decomposition analysis on Qwen-7B attributes the gap to information overload and gradient dilution from low-density tokens, and a tail-truncation budget sweep confirms a sharp phase transition between compact and verbose inputs.
These results demonstrate that for trace-based verdict prediction, learning efficiency correlates far more strongly with information density than with information quantity.
\end{abstract}

\input{introduction}
\input{method}
\input{experiments}
\section{Discussion}

The central finding is that for trace-based verdict prediction, learning efficiency strongly correlates with information density rather than information quantity.
Exception events concentrate the runtime's primary fault signal---error type, stack frame, and message---in a compact format (median 150 tokens), while non-exception components (variable states, loop iterations, control flow) record \emph{how} the program ran, not \emph{whether} it failed.

This density perspective explains three key results: (1)~adding trace information beyond exceptions provides no benefit (Table~\ref{tab:ablation}), because low-density content competes for attention; (2)~a tail-truncation budget sweep reveals a sharp phase transition at ${\sim}$500 tokens (Table~\ref{tab:tail_cross_model}), matching the boundary where diagnostic events dominate; (3)~the SFT advantage over regex concentrates on caught-exception boundary cases, confirming semantic signal extraction (\S\ref{sec:ablation}).
The narrow cross-model gap (Table~\ref{tab:cross_model}) supports that signal density, not model capacity, is the binding constraint.

\paragraph{Why exception\_only outperforms ae\_only.}
Although \texttt{ae\_only} is a strict superset, it underperforms \texttt{exception\_only} (Table~\ref{tab:ablation}), indicating extra content introduces noise rather than signal.

\paragraph{SFT vs.\ regex baselines.}
A natural question is whether SFT is necessary when exception signals are so discriminative.
We evaluate two regex baselines: detecting \texttt{[ERROR]} tokens ($B_{\text{error}}$, 93.45\%)
and detecting the \texttt{[No exception data in trace]} marker ($B_{\text{marker}}$, 99.33\%).
The optimal regex still misclassifies 8 samples, while SFT misclassifies only 2
(McNemar $p{=}0.014$). All 6 discordant samples involve \emph{caught exceptions}---programs
where a \texttt{try/except} block catches a \texttt{ValueError} but execution succeeds.
SFT learns to distinguish caught from uncaught exceptions; regex cannot.
Beyond accuracy, SFT provides calibrated probability estimates (ECE${=}$0.003),
format robustness (96.55\% without markers; Appendix~\ref{sec:nomarker}), and requires no task-specific engineering.

\paragraph{Why zero-shot and few-shot prompting fall short.}
Zero-shot prompting spans a 42pp range across models (\S\ref{sec:additional}), which SFT compresses to sub-1pp; few-shot ICL improves marginally, confirming the bottleneck is verdict calibration rather than exception recognition.


\paragraph{Near-perfect floor, not hard ceiling.}
The 2 errors (both \texttt{GeneratorExit}; Appendix~\ref{sec:error_overlap}) are shared across configurations; CodeLlama-13B seed 123 achieves 100\%, confirming 99.83\% is a convergence \emph{floor}, signal-limited rather than capacity-limited.

\paragraph{Tokenizer efficiency.}
CodeLlama's tokenizer produces ${\sim}$20\% more tokens than Qwen's on identical inputs, and a label-only loss variant collapses at tail-2048, indicating LoRA long-sequence instability (Appendix~\ref{sec:tokenizer_analysis}).

\paragraph{Tier hierarchy and methodological implications.}
The two-tier structure (Table~\ref{tab:cross_model}) suggests a fallback design: extract exception events first, resort to Tier~2 only when no exception is raised.
Loop benefit varies by architecture (Qwen-7B: +10.9pp; CodeLlama-13B: near-zero), likely reflecting differences in pre-training coverage of loop patterns.
The methodological insight---that information density strongly correlates with learning efficiency---transfers to any trace-based task where compact diagnostic events exist within verbose traces.

\paragraph{Density and cross-model robustness.}
Tier~1 and Tier~2 exhibit contrasting cross-model stability (Table~\ref{tab:cross_model}): \texttt{exception\_only}'s range is ${<}$1pp, while \texttt{loop\_only} spans ${\sim}$9pp.
However, $N{=}3$ models is insufficient to establish a continuous relationship, and Tier~1 signals are both high-density \emph{and} direct fault indicators, so density cannot be disentangled from signal type.
Applying label-only loss to DeepSeek full traces (91.7\% truncated) causes complete collapse (28.5\%), confirming that trace content density, not loss function, is the binding constraint.


\paragraph{Caveats on the density decomposition.}
Our 24\%/76\% gradient dilution vs.\ information overload decomposition
has two important caveats. First, the two factors are not orthogonal:
\texttt{label\_only} simultaneously changes the gradient source distribution
and the effective learning rate, so the decomposition is an approximation
rather than a causal isolation. Second, the decomposition is verified only
on the Qwen family (two scales); CodeLlama and DeepSeek \texttt{label\_only}
conditions collapse entirely, preventing cross-family replication.
Sequence length is an additional confounding variable that our design does not isolate.

\paragraph{Exception events capture test verdict signals, not just crashes.}
Although 55.5\% of failures are silent logic errors, \texttt{exception\_only} achieves 100\% fail recall because test assertions generate \texttt{AssertionError} events, confirming ``exception events'' encompasses both crashes and assertion failures.

\begin{figure}[t]
\centering
\includegraphics[width=\columnwidth]{training_dynamics}
\caption{Training dynamics comparison. Exception-only (blue) converges to 99.33\% within 50 steps and saturates at 99.83\% by step 150. Full-trace with label-only loss (red, dashed) oscillates between 82\% and 89\% across all 444 training steps without stable convergence, consistent with overfitting to spurious patterns in verbose trace content.}
\label{fig:training_dynamics}
\end{figure}

\paragraph{Information density and learning dynamics.}
Figure~\ref{fig:training_dynamics} shows the contrast: \texttt{exception\_only} converges rapidly, while \texttt{full\_trace} with label-only loss oscillates throughout---a signature of overfitting to spurious patterns (Appendix~\ref{sec:training_dynamics}).
Seed variance negatively correlates with input density (Table~\ref{tab:ablation}).

\section{Related Work}

\paragraph{Execution Traces for Code LLMs.}
TracePile \citep{chen2025tracepile} scaled trace-augmented training to 2.6M samples; \citet{wang2025dcsh} found trace utility task- and format-dependent, which we extend by identifying exception events as the dominant signal.
NExT \citep{ni2024next} and ET-CoT \citep{eguchi2025etcot} used trace chain-of-thought for repair and execution reasoning; \citet{ferrando2025execute} found compact traces comparable to full traces, supporting our density finding.
\citet{liu2023codeexecution} pre-trained on execution traces; CRUXEval \citep{gu2024cruxeval}, extended to 19 languages \citep{xu2025cruxevalx}, established execution reasoning benchmarks.
None decompose which trace components contribute to \emph{classification} accuracy.

\paragraph{Code Verdict and Patch Classification.}
TRACED \citep{ding2024traced} showed execution-aware pretraining improves vulnerability detection; DeepDFA \citep{steenhoek2024deepdfa} and learned program slicing \citep{yadavally2024slicing} combined classical analysis with deep learning.
In patch correctness, \citet{tian2020patchcorrectness}, CC2Vec \citep{hoang2020cc2vec}, LLM4PatchCorrect \citep{zhou2024llm4patch} (84.4\%), and APPT \citep{appt2024} (79.7\%) developed representation learning of code changes.
Verdict prediction is an instance of the longstanding test oracle problem \citep{barr2015oracle}; FuseFL \citep{widyasari2024fusefl} and \citet{foster2025mutation} demonstrated structured fault signals; the equivalent-mutant problem \citep{papadakis2019mutation} poses an analogous classification challenge.
None isolate individual trace components or compare against rule-based baselines.

\paragraph{Fault Localization.}
Spectrum-based fault localization (SBFL) ranks program elements by suspiciousness from test spectra \citep{abreu2007accuracy,jones2002visualization,wong2014dstar}; recent LLM-based approaches \citep{kang2024autofl,qin2024agentfl} share our insight that anomalous execution events carry the strongest diagnostic signal, with verdict prediction serving as a prerequisite to FL.

\paragraph{Trace Compression and Selection.}
TraceCoder \citep{huang2026tracecoder} and InspectCoder \citep{wang2025inspectcoder} use selective trace extraction; RepairAgent \citep{bouzenia2025repairagent}, Agentless \citep{xia2024agentless}, SWE-agent \citep{yang2024sweagent}, and OpenHands \citep{wang2025openhands} incorporate execution feedback in repair workflows \citep{jimenez2024swebench}, yet which trace components carry verdict signal most efficiently remains unstudied.
LLMLingua \citep{jiang2023llmlingua} and Selective Context \citep{li2023selectivecontext} showed that pruning low-information tokens preserves accuracy; \citet{haque2025traces} found raw traces provide limited repair benefit.
Our decomposition (${\sim}$24\% gradient dilution vs.\ ${\sim}$76\% information overload) provides the first quantitative framework for \emph{why} compression helps.

\paragraph{Calibration and Defect Prediction.}
Neural network calibration \citep{guo2017calibration,kadavath2022language,spiess2025calibration} studies confidence-correctness alignment; our ECE ${<}$ 0.007 suggests compact inputs improve calibration. Just-in-time defect prediction \citep{kamei2013large,hoang2019deepjit} shares our fault detection motivation but operates on code change features.

\paragraph{Training Efficiency and Data Selection.}
Curriculum learning \citep{bengio2009curriculum} and data selection methods \citep{swayamdipta2020cartography,sachdeva2024dataefficient} show that data quality improves convergence, consistent with our finding that compact exceptions converge rapidly while verbose traces oscillate.

\input{conclusion}

\section*{Limitations}

Our study has several limitations.
First, our primary experiments use a single dataset (591 problems from HumanEval and MBPP-sanitized); while a cross-dataset evaluation on DebugBench (\S\ref{sec:debugbench}) confirms that exception dominance generalizes to synthetically injected bugs, generalization to other languages, larger codebases, or real-world fault distributions remains untested.
Second, while we evaluate four models spanning 1.5B to 13B parameters across three architecture families, all are code-specialized; generalization to general-purpose models or larger scales remains untested.
Third, DeepSeek VOTA seed 42 (84.45\%) was excluded as a training failure due to tokenizer-driven truncation (Appendix~\ref{sec:tokenizer_analysis}); reported VOTA means use the remaining 2 seeds.
Qwen \texttt{ae\_only} exhibits mild seed variance ($\sigma{=}$0.40pp; seed 123 = 99.08\% vs.\ seeds 42/456 = 99.92\%); other conditions use a single seed.
Fourth, the initial DeepSeek \texttt{full\_trace} run was invalidated by context overflow; a re-run with extended context yields 79.16\% (CI [76.80\%, 81.30\%]), 2.60pp below Qwen (Table~\ref{tab:cross_model}).
Fifth, the candidate solutions were generated by Qwen-2.5-Coder-7B-Instruct, which also serves as one of our evaluation models; this does not affect verdict prediction but may influence the pass/fail distribution.
Sixth, self-familiarity bias may inflate absolute accuracy: both Qwen-2.5-Coder and DeepSeek-Coder were likely exposed to HumanEval and MBPP solutions during pretraining, so the models may partially recognize solutions rather than purely reasoning from trace signals; our cross-dataset evaluation on DebugBench mitigates but does not eliminate this concern.
Seventh, the error distribution in our test set may not reflect real-world fault distributions: exceptions (especially \texttt{TypeError}, \texttt{ValueError}) dominate LLM-generated failures, but production code may exhibit a broader spectrum of silent logic errors where no exception is raised, potentially reducing the effectiveness of exception-based verdict prediction.
Finally, our analysis is limited to verdict prediction; whether the information density advantage extends to other trace-based tasks is an open question.

Our primary experiments consist entirely of LLM-generated Python short functions. Generalization to human-written code,
other programming languages, and larger multi-file projects remains unverified.
Format invariance was verified only on Qwen-7B \texttt{exception\_only}; cross-model and cross-condition
format sensitivity remains untested.

\section*{Ethics Statement}

This work uses publicly available benchmarks (HumanEval, MBPP-sanitized, DebugBench) and does not involve human subjects, crowdsourced annotations, or personally identifiable information.
All models evaluated are open-source and publicly accessible.
Training uses parameter-efficient LoRA fine-tuning on two RTX PRO 6000 Blackwell GPUs, with each experiment using a single GPU.
Total compute is approximately 320 GPU-hours across all conditions, models, seeds, and iterative refinements (including approximately 170 hours for successful experiments; individual runs range from 0.5h for small models to 2.2h for CodeLlama-13B, with the total reflecting 4 models $\times$ multiple conditions $\times$ multiple seeds $\times$ iterative development).
We do not foresee direct negative societal impacts from this work; verdict prediction is a developer-facing tool for automated testing.

\bibliography{references}

\appendix
\input{appendix}

\end{document}
